\chapter{Untyped \lCalc}

In the early 20th century, mathematicians such as Bertrand Russell, David Hilbert, and Kurt Gödel were trying to set complete, consistent, and decidable foundations for mathematics. Within this context, Alonzo Church, in his effort to formalize the notion of computability, introduced a minimal symbolic language based upon function abstraction and function application, this system is now called \lcalc.

In 1936, Church used lambda calculus to address Hilbert’s Entscheidungsproblem --whether a mechanical method could determine the truth of any first-order logic statement. He proved no such algorithm exists, establishing the undecidability of first-order logic.

One of the nuances formalized by the \lcalc \ is the distinction between extensional and intensional equality. The extensional approach to equivalence states that two functions are equivalent if they share input output pairs for every possible input. Its intensional counterpart extends this notion of equality by requiring that the procedures that compute these pairs share complexity i.e. they take the same steps toward yielding a result.\footnote{
  Let \( p \) be a sufficiently large prime, and let \( f, g : \mathbb{Z}_p \to \mathbb{Z}_p \) be defined by \( f(x) = x^2 \) and \( g(x) = \log_a(a^{x+2}) \), where \( a \in \mathbb{Z}_p^\times \) is a fixed primitive root. Although \( f \) and \( g \) are extensionally equal, i.e., they yield the same output for all inputs in \( \mathbb{Z}_p \), they are intensionally distinct. The function \( f \) performs a simple squaring operation, while \( g \) requires evaluating a discrete logarithm, intractable in general \cite{bonilla2021}.}

The Lambda Calculus together with Type Theory lay the theoretical groundwork that shaped the design of languages like LISP, Haskell and Erlang to name a few. More advanced type systems than those present in general purpose programming languages enable computers to verify mathematical statements and produce rigorous proofs.

\subsubsection{\centering The Intuitive Notion}

In order to get acquainted with the \lcalc, let us develop a simple example to familiarize ourselves before we begin with a more formal approach to this discipline. Consider the function $f(x) = x + 1$, the most straightforward way to express this using \lcalc \ would be $(\la x . x + 1 )$, where the lambda denotes that $x$ is being captured and used as a parameter to perform some computation. Evaluating $f(2)$ using this newly created lambda term would look like this: $(\la x . x + 1)(2) \mathrel{\leadsto} (2 + 1) \mathrel{\leadsto} 3$. This process of reducing a $\la$-expression is referred to as \bred, it will be covered formally later in the chapter.

\section{\centering The Language of \lCalc}

\begin{definition} The set of all lambda terms \( \La \) is defined inductively as follows:
  \label{def:lambda-terms-1}
  \begin{itemize}
  \item (Variable) If \( x \in \Vset \), then \( x \in \La \).
  \item (Abstraction) If \( x \in \Vset \) and \( M \in \La \), then \( (\la x. M) \in \La \).
  \item (Application) If \( M, N \in \La \), then \((M N) \in \La \).
  \end{itemize}
  Where $\Vset = \{x, y, z, ... \}$ represents a countably infinite set of variable names. These notions are conveniently synthesized into the following abstract syntax:
  \[
    \La \, := \, \Vset \, \mid \, \la x . \, \La \, \mid \, \La \, \La
  \]
  \noindent
  that directly translates into the following Haskell snippet:
  \begin{minted}[fontsize=\small]{haskell}
    data Term = Var String
          | Abs String Term
          | App Term Term
  \end{minted}
\end{definition}

The key takeaway of this definition is that abstraction and application when paired with a reduction procedure, capture the meaning of function in a way that is abstract yet useful, and enables computation.

\begin{notation}
  We use uppercase letters, e.g. M, N, L, P, Q to denote elements of $\La$ and lowercase e.g. x, y, z to denote elements of $\Vset$.
\end{notation}

\begin{remark}
  From the definition above we extract that \( \lambda x . x + 1 \notin \La \), given this circumstance, the need to find an accurate representation for numbers within this definition arises. See Definition \ref{def:church-naturals}
\end{remark}

\begin{example} \label{ex:lambda-terms} Some examples of valid $\la$ terms generated using the syntax from \ref{def:lambda-terms-1}:
  \( y \),
  \( (\la x. (x x)) \),
  \( (\la x. (\la y. x)) \),
  \( (((\la x. (x y)) (\la y. y)) z) \).
  The Haskell counterpart for the second term is:
  \begin{minted}[fontsize=\small]{haskell}
    Abs "x" (Abs "y" (Var "x"))
  \end{minted}
\end{example}
As a convention to avoid notational cluttering, the outermost parenthesis can be omitted e.g. $\la x.x $ is read as $ ( \la x.x ) $. Also, application is left-associative  and binds tighter than abstraction, so $M N P$ is parsed as $(M N) P$ and $\la x . M N$ means $\la x . (M N)$, not $(\la x . M) N$. Likewise, abstraction is right-associative e.g. $\la x . \la y . M$ is $\la x . (\la y . M)$ and binds more weakly than application. Haskell Curry introduced what we now refer to as currying, its main purpose is to avoid unnecessary lambdas. Below is a summary of common notational conventions used throughout the text:
\begin{itemize}
\item \( \la x . \, x \iff ( \la x . \, x ) \)
\item \( M \, N \, P \iff (M \, N) \, P \)
\item \( \la x. M \, N \iff \la x. (M \, N) \)
\item \( \la x. \la y. f(x,y) \iff \la xy. f(x,y) \)
\end{itemize}
All these are just to keep notation straight-forward, since the formal definition does not leave precedence under-specified thanks to parenthesis. Those interested in the implementation need not worry, as Haskell enforces an explicit order on constructors.
\section{\centering Reduction of terms}
Now that the \lcalc \ language has been introduced, it is time we move on to the mechanics of computation, manifested through $\beta$-equivalence. $\alpha$ and $\eta$ equivalence will also be introduced. These two are more difficult to grasp. Informally, it could be said that $\alpha$-equivalence is equivalence programs up to renaming of local variables. In this same scenario, $\eta$-equivalence would singnify equivalence of programs up to wrapping or unwrapping them with an identity function.
\subsubsection{\centering Subterms}
From the ground up, we define what a subterm is, since the definition of the language is self-referential and terms might contain other terms.
\begin{definition} Given $T \in \La$, $\functionfont{Sub} : \La \to \mathcal{P}(\Vset)$ maps a term to the set of it's subterms.
  \begin{align*}
    &\Sub x = \{x\} \\
    &\Sub {MN} = \Sub M  \cup \Sub N \cup \{MN\} \\
    &\Sub {\la x . M} = \Sub M  \cup \{\la x. M\}
  \end{align*}
  Where $x \in \Vset$, $M, N \in \La $ and $\mathcal{P}$ denotes the power set. In Haskell:
  \begin{minted}[fontsize=\small]{haskell}
    subTerms :: Term -> Set.Set Term
    subTerms (Var x) = Set.singleton (Var x)
    subTerms (App m n) =
        Set.insert (App m n) (Set.union (subTerms m) (subTerms n))
    subTerms (Abs x m) =
        Set.insert (Abs x m) (subTerms m)
  \end{minted}  
\end{definition}
\begin{remark}
  Let \( \preceq \) be the subterm relation on \lterms, defined by
  \[
    M \preceq N \iff M \in \Sub{ N }
  \]
  \( \preceq \) is a partial order relation. This statement is follows intuitively from the definition, it can be proven using induction on the derivation tree.
\end{remark}
\begin{definition} A term $ S \in \Sub M, \ M \in \La$ is said to be proper if $S \neq M$. $\Sub{M} \setminus \{M\}$ would be the set of proper subterms.
\end{definition}
\begin{definition} Let \( M \equiv \lambda x. N \) be a \lterm:
  \begin{itemize}
  \item The variable \( x \) is the binding variable of the abstraction \( \lambda x. N \).
  \item An occurrence of a variable \( x \) in \( M \) is bound if \( x \in \Sub { \lambda x. N } \).
  \item Free variables are those that are not bound.
  \end{itemize}
\end{definition}
\begin{example}
  Take $ M \equiv \la x . \, x \, y $, in this example, the $x$ in $\la x. $ is binding, $y$ is free, and $x$ is bound. This means that $y$ is used literally, and that $x$, as denoted by $\la x.$ is abstracted and thus subject to being replaced within the context of said abstraction. A subexample might make this explicit: $ (\la x . \, x \, y) \, M \curly M \, y $. The $y$ remained while the $x$ was used as a placeholder for $M$.
\end{example}
\begin{note}
  This taxonomy is key to understanding concepts such as \aequiv \ intuitively.
\end{note}
\subsubsection{\centering Free Variables}
The set of Free Variables $\functionfont{FV}$, those that will not get replaced during the process of reduction, is computed using:
\begin{definition} Given $T \in \La $, $\functionfont{FV} : \La \to \mathcal{P}(\Vset) $ outputs the set of free variables for $T$:
  \begin{align*}
    & \FV x = \{x\}\\
    & \FV {MN} = \FV M \cup \FV N \\
    & \FV {\lambda x . M} = \FV M \setminus \{x\}
  \end{align*}
  Where $ x \in \Vset $ and $ M, N \in \La $. In Haskell:
  \begin{minted}[fontsize=\small]{haskell}
    freeVars :: Term -> Set.Set String
    freeVars (Var x) = Set.singleton x
    freeVars (App m n) = Set.union (freeVars m) (freeVars m)
    freeVars (Abs x m) = Set.delete x (freeVars m)
  \end{minted}  
\end{definition}
\begin{example} Computing the free variables of $\lambda x . \lambda y . y x z$. The term well-formed, so we proceed:
  \begin{align*}
    \FV {\lambda x . \lambda y . y x z} &= \FV {\lambda y . yxz} \setminus \{x\} \\
                                        &= \FV {yxz} \setminus \{x, y\} \\
                                        &= \{x, y, z\} \setminus \{x, y\} \\
                                        &= \{z\}
  \end{align*}
  Of course, $z$ is the only free variable in the expression, since it is the only variable that is not captured by some $\la$-abstraction.
\end{example}
\begin{definition}
  A term $M \in \La$ is closed if and only if $\FV M = \emptyset$, closed lambda terms are often referred to as combinators. $\La^0$ is the set of all closed lambda terms.
\end{definition}
\begin{note} They are called combinators since having no free variables means that all variables are bound, and thus, a closed term just combines bound variables.
\end{note}
\subsubsection{\centering Variable Substitution}
We approach the definition for \aequiv \, to this end we have to introduce what it means to substitute a variable.
\begin{definition} Capture-Avoiding Substitution Rules
  Let \( M, N \in \La \) and \( x, y, z \in \Vset \). The substitution \( M[x := N] \) is defined inductively as:
  \[
    \begin{aligned}
      x[x := N]                       & = N \\
      y[x := N]                       & = y && y \neq x \\
      (M_1\, M_2)[x := N]             & = (M_1[x := N])\,(M_2[x := N]) \\
      (\lambda x. M)[x := N]          & = \lambda x. M \\
      (\lambda y. M)[x := N]          & = \lambda y. (M[x := N]) && y \neq x, y \notin \FV { N } \\
      (\lambda y. M)[x := N]          & = \lambda z. (M[y := z][x := N]) && y \neq x, y \in \FV { N }
    \end{aligned}
  \]
  where \( z \notin \FV { M } \cup \FV { N } \). Implementing the latter in Haskell requires a bit more code than previous definitions, along with making some practical tradeoffs:
  \begin{minted}[fontsize=\small]{Haskell}
  subst :: String -> Term -> Term -> Term
  subst x s v@(Var y)
      | x == y = s
      | otherwise = v
  subst x s l@(Abs y body)
      | x == y = l
      | Set.notMember y (freeVars s) = Abs y (subst x s body)
      | otherwise =
          let used = Set.unions [allVars body, allVars s, Set.singleton x]
              z = freshVar used
              body' = rename y z body
          in Abs z (subst x s body')
  subst x s (App a b) = App (subst x s a) (subst x s b)
\end{minted}      
      \begin{note}
        \texttt{freshVar} generates a new name for binding variables such that \( z \notin \FV { M } \cup \FV { N } \). More details are available in the source code.
      \end{note}
  \end{definition}
Hopefully the meaning of the first three rules is clear, they ancapsulate renaming in the context of variables and applications. The subsequent rules are about abstractions, rule four states that renaming a binding variable has no effect, renaming only modifies a term whenever the renamed variable is different from the binding variable. Rule five is the general case, that is, when the substitution just moves in together with the abstracted term. Then there is the last case, when the binding variable is a bound occurrence in the term we are substituting, which poses a problem, to see why, let us work through an example:
\begin{example} If we mistakenly substitute $x$ by \( (\la w . \, w \, y) \) in \( (\la y . \, y \, x ) \) using rule five, we get:
  \[
    (\la y . \, y x) [x := (\la w . \, w \, y)] = (\la y . \, y \, (\la w. \, w \, y))
  \]
  This is problematic since the free variable $y$ became bound after the rename. Free and bound variables should maintain their freedom or boundedness after a rename. The correct way to substitute avoiding the name clash:
  \[
    (\la y . \, y x) [x := (\la w . \, w \, y)] = (\la z . \, z \, (\la w. \, w \, y))
  \]
  Binding variable serve as placeholders, therefore, renaming should not affect the integrity of an expression.
\end{example}

\begin{note}
  Hence the use of Capture Avoiding alongside Substitution.
\end{note}


One notable result regarding substitution is the substitution lemma:

\begin{lemma} Let \(M, N, P\) be lambda terms and \(x, y\) be variables. Then:
\[
M[x := N][y := P] =
\begin{cases}
M[x := N[y := P]], & \text{if } x = y, \\
M[y := P][x := N[y := P]], & \text{if } x \neq y \text{ and } y \notin \mathrm{FV}(N).
\end{cases}
\]
\end{lemma}

\begin{proof}[Proof Sketch]
  By structural induction on $M$, using $\alpha$-conversion to avoid capture and
assuming bound variables $x,y,N,P$ are fresh.

\begin{itemize}
\item Base case: Direct from the definition of substitution; check $z = x$,
$z = y$, and $z$ distinct from both.
\item Inductive step on applications, where substitution distributes:
\[
(M_1M_2)[x:=N] = M_1[x:=N]\,M_2[x:=N],
\]
so both statements follow immediately from the induction hypothesis on $M_1$ and $M_2$.
\item Inductive step on abstractions with $z$ fresh,
\[
(\lambda z.M_0)[x:=N] = \lambda z.(M_0[x:=N]),
\]
and the induction hypothesis applies to $M_0$.
\end{itemize}
\end{proof}

In the previous example a bound variable is renamed to avoid a naming collision and avoid binding a previously free variable, \aequiv \ is the generalization of this technique, it defines equivalence up to renaming of binding variables.

\begin{notation}
  In many of the upcoming definitions inference rules are used, here is a simple introduction for those unfamiliar with the concept:
  \[
    \inferrule*[right=\textsc{Name}]{
      P_1 \quad P_2 \quad \cdots \quad P_n
    }{
      C
    }
  \]
  where $P_i$ are judgments assumed to hold, $C$ is the judgment inferred from the premises and \textsc{Name} is an optional identification label.
\end{notation}
\subsubsection{\centering $\alpha$-equivalence}
\begin{definition} Two lambda terms \( M, N \in \La \) are \aequivlt, written \( M =_\alpha N \), if they are structurally identical except for the names of bound variables.
  Formally:
  \[
    \begin{array}{@{}l@{\quad}l@{}}
      \inferrule*[right=\textsc{Alpha}]
      { }
      { \lam{x}{L} =_\alpha \la y . \, L \subst{x}{y} }
      &
        \inferrule*[right=\textsc{Abs}]
        { L =_\alpha L' }
        { \lam{x}{L} =_\alpha \lam{x}{L'} } \\[2ex]
      \inferrule*[right=\textsc{AppL}]
      { L_1 =_\alpha L_1' }
      { L_1\,L_2 =_\alpha L_1'\,L_2 }
      &
        \inferrule*[right=\textsc{AppR}]
        { L_2 =_\alpha L_2' }
        { L_1\,L_2 =_\alpha L_1\,L_2' }
    \end{array}
  \]
  $\alpha$-equivalence is an equivalence relation:
  \[
    \begin{aligned}
      &\text{(Reflexivity)} \quad && M =_\alpha M \\
      &\text{(Transitivity)} \quad && M =_\alpha N \text{ and } N =_\alpha P \Rightarrow M =_\alpha P \\
      &\text{(Symetry)} \quad && M =_\alpha N \Rightarrow N =_\alpha M
    \end{aligned}
  \]
\end{definition}
\begin{example} A simple equivalence class and two samples of not $\alpha$-equivalent \ pairs.
  \[
    \la x.\la y. \, y =_\alpha \la a.\la b. \, b =_\alpha \la z.\la x. \, x
  \]
  \[
    \la x. \la y. \, x \not=_\alpha \la x. \la y. \, y \quad \la x. \, x \, y \not=_\alpha \la x. \, x \, z 
  \]
  Remember terms are $\alpha$-equivalent \ only when bounded names are properly substituted.
\end{example}
\begin{remark}
  \aequiv \ comes naturally once substitution is defined, since the names of bound variables are just placeholders that will eventually hold an argument.
\end{remark}

\begin{note}
  Variable naming is a nightmare. De Bruijn indexing is a way of methodically assigning numbers to variable names, indicating how many binders away a variables binder is. This system eliminates the need for $\alpha$ equivalence to exist.
  \[
    \la x.\, x \Rightarrow \la.\, 0 \quad \la x. \la y.\, x \Rightarrow \la. \la.\, 1 \quad \la x. \la y.\, y \Rightarrow \la. \la. \, 0\,
  \]
  We are better at names than we are at numbers, which is why we will prefer using symbols over De Brujn indexing. Programs like automated theorem provers and compilers do use similar naming strategies to simplify their work.
\end{note}
\subsubsection{\centering $\beta$-equivalence}
Defines the means to perform computation in the \lcalc, applying substitution as a means of reduction for expressions. This is similar to how one simplifies an arithmetic expression, only it is more powerful, since arithmetic expressions lack expressiveness. It fails to reproduce branching for instance.
\begin{definition} Single-step $\beta$-reduction $\step$ is defined using the following inference rules:
  \[
    \begin{array}{@{}l@{\quad}l@{}}
      \inferrule*[right=\textsc{Beta}]
      { }
      { (\lam{x}{M})\,N \step M\subst{x}{N} }
      &
        \inferrule*[right=\textsc{Abs}]
        { M \step M' }
        { \lam{x}{M} \step \lam{x}{M'} } \\[2ex]
      \inferrule*[right=\textsc{AppL}]
      { M_1 \step M_1' }
      { M_1\,M_2 \step M_1'\,M_2 }
      &
        \inferrule*[right=\textsc{AppR}]
        { M_2 \step M_2' }
        { M_1\,M_2 \step M_1\,M_2' }
    \end{array}
  \]
  where \( M, N \in \La \). In Haskell:
  \begin{minted}[fontsize=\small]{Haskell}
  step :: Term -> Maybe Term
  step (App (Abs x m) n) = Just (subst x n m)
  step (App a b) = case step a of
      Just a' -> Just (App a' b)
      Nothing -> case step b of
          Just b' -> Just (App a b')
          Nothing -> Nothing
  step (Abs x t) = case step t of
      Just t' -> Just (Abs x t')
      Nothing -> Nothing
  step (Var _) = Nothing
  \end{minted}
\end{definition} 
Single \bred s can be applied succesively, which induces the following definition
\begin{definition}
  \label{def:multi-step-beta-reduction}
  Multi-step \bred. We write $M \step^* N$ if and only if there exists a sequence of one-step $\beta$-reductions starting from $M$ and ending at $N$.\footnote{The convention in the literature is to use $\rightarrow_\beta$ for \bred s $\twoheadrightarrow_\beta$ for multi-step \bred s. We prefer to use $\curly$ and $\curly^n$ since this allows us to numerate \bred \ steps e.g. $ (\la f. \la x . \, x) \, M \, N \curly^2 $ N. Instead of $\curly^n$, we use $\curly^*$ whenever the number of steps is undefined.}
  \label{def:zero-more-bred}
  \[
    M \equiv M_0 \step M_1 \step M_2 \step \cdots \step M_{n-1} \step M_n \equiv N
  \]
  That is, there exists an integer $n \geq 0$ and a sequence of terms $M_0, M_1, \dots, M_n$ such that: $ M_0 \equiv M $ and $ M_n \equiv N $ for all $ i $ with $ 0 \leq i < n $, we have $ M_i \step M_{i+1} $.
\end{definition}
The process of succesively $\beta$-reducing a term leads raises the question of whether all terms can be reduced indefinitly. We know the answer to this, since for instance, variables in $\Vset$ cannot be reduced. The subsequent question is then, Do all all terms have some kind of normal form that cannot be reduced further? This is answered in Section \ref{sec:recursion-fixed-points}. We now introduce a formal definition of normal form that will later allow us to gain insight into the normalization properties of the \lcalc.

\begin{definition} A \(\lambda\)-term \(M\) is in normal form if there is no term \(N\) such that \( M \curly N \). Equivalently, \(M\) contains no subterm of the form \( (\la x . \, P) \, Q \) that can be reduced by beta-reduction. We say a term \(M\) is irreducible under beta-reduction.
\end{definition}

\begin{definition}
  A reducible expression, or redex for short, is any subterm of the form 
  \( (\la x . \, P) \, Q \). Such a term can be reduced by a single beta-reduction step, resulting in \( P[x := Q] \).
\end{definition}

\begin{definition} \bequiv \ relation, equivalence of \lterms \ through \bred s.
  \[
    M =_\beta N \iff M \step^* N \lor N \step^* M
  \]
\end{definition}
\begin{note}
  $=_\beta$ is the symetric closure of $\curly^*$.
\end{note}
We now explore a key aspect of the \lcalc, in the jargon referred to as first-class citizenship, or high-order, a property of terms that restrains them from being exclusively classified as either function or argument. The example below sheds some light on the matter:
\begin{example} High Order by example. Let us $\beta$-reduce the \lterm \ $ A \, B \, C $ and inspect the consequences:
  \label{ex:high-order-by-example}
  \begin{align*}
    \overbrace { (\lambda f . \lambda y . f f y) }^{A} \overbrace{ (\lambda x . x + 1) }^{B} \overbrace{(2)}^{C}
    &\curly (\lambda y .(\lambda x . x + 1) (\lambda x . x + 1)y)(2) \\
    &\curly (\lambda x . x + 1)(\lambda x . x + 1)(2) \\
    &\curly^* (\lambda x . x + 1) (3) \\
    &\curly^* (4)
  \end{align*}
  Taking a closer look, $A$, applies $f$ twice to $x$, as a consequence, $B$ is being applied twice to $C$. And so, very easily, we have implemented the $ x + 2 $ function via a sequential application of $x + 1$ to $C$.
\end{example}
\begin{remark}
  As seen above, we can pass what we understand as conventional functions as arguments since they are treated alike.  One can intuitively apreciate the computational expresivenes this brings with it, and how this syntactic-semantic homogeneity sets the \lcalc \ apart from the classical set-theoretic approach to functions.
\end{remark}

\subsubsection{\centering $\eta$-equivalence}
This last notion of equivalence, refers to a program being equivalent to itself wrapped with an identity function:
\begin{definition} Two lambda terms \( M \) and \( N \) are said to be $\eta$-equivalent, written \( M =_\eta N \), if:
  \[
    \begin{array}{@{}l@{\quad}l@{}}
      \inferrule*[right=\textsc{Eta}]
      { x \notin \FV M }
      { (\lam{x}{M \, x}) =_\eta M   }
      &
        \inferrule*[right=\textsc{Abs}]
        { M =_\eta M' }
        { \lam{x}{M} =_\eta \lam{x}{M'} } \\[2ex]
      \inferrule*[right=\textsc{AppL}]
      { M_1 =_\eta M_1' }
      { M_1\,M_2 =_\eta M_1'\,M_2 }
      &
        \inferrule*[right=\textsc{AppR}]
        { M_2 =_\eta M_2' }
        { M_1\,M_2 =_\eta M_1\,M_2' }
    \end{array}
  \]

  The abstraction \( \lambda x.\, (M\;x) \) performs the same operation as \( M \), so if \( x \) is not used freely within \( M \), the abstraction serves no purpose. This is shown in \textsc{Eta}, the rest of the inference rules propagate this equivalence throughout the whole term.
\end{definition}
\begin{note}
  Two terms are observationally equivalent if they cannot be distinguished by any context that can be applied to them. $\eta$-equivalence can be understood as a form of weak extensionality\footnote{Understood as we presented it in the introduction}, since it does not fully encode observational equivalence.
\end{note}

\subsubsection{\centering $\beta\eta$-reduction}

To later prove the Church--Rosser property, we introduce a set of rules that captures reduction taking extensionality into account, we combine both $\beta$ and $\eta$-reduction into a single relation.

\begin{definition}
We write $M \betared N$ if $M$ reduces to $N$ by one application
of either a $\beta$ or an $\eta$-rule, possibly in a subterm. The rules are:

\[
  \begin{array}{@{}l@{\quad}l@{}}
    \inferrule*[right=\textsc{Beta}]
    { }
    { (\lambda x. M)\,N \;\betared\; M[x := N] }
    &
      \inferrule*[right=\textsc{Eta}]
      { x \notin \FV{M} }
      { \lambda x.(M\,x) \;\betared\; M } \\[2ex]
    \inferrule*[right=\textsc{AppL}]
    { M_1 \betared M_1' }
    { M_1\,M_2 \;\betared\; M_1'\,M_2 }
    &
      \inferrule*[right=\textsc{AppR}]
      { M_2 \betared M_2' }
      { M_1\,M_2 \;\betared\; M_1\,M_2' } \\[2ex]
  \end{array}
\]
\[
  \inferrule*[right=\textsc{Abs}]
  { M \betared M' }
  { \lambda x.M \;\betared\; \lambda x.M' }
\]

\end{definition}

\begin{definition}
  We write $M \betareds N$ if $M$ reduces to $N$ by zero or more single-step $\beta\eta$-reductions, the same way we did for $\beta$-reduction in  \ref{def:multi-step-beta-reduction}.
\end{definition}


